\section{Hamiltonian Ising sebagai Representasi Sistem Kuantum}

\subsection{Sejarah dan Evolusi Model Ising}
Model Ising awalnya dikembangkan oleh Ernst Ising dan Wilhelm Lenz pada awal abad ke-20 sebagai instrumen teoretis untuk mendeskripsikan fenomena feromagnetisme dalam kerangka kerja mekanika statistik \cite{ising_1925}. Fokus utama dari model ini adalah untuk memodelkan transisi fase pada material magnetik melalui mekanisme interaksi sederhana antara partikel-partikel diskrit yang saling bertetangga di dalam suatu sistem fisik. Interaksi mikroskopis tersebut digambarkan melalui parameter $J_{ij}$ yang menyatakan kekuatan kopling antar-\textit{spin} dan variabel $s_i \in \{+1, -1\}$ sebagai orientasi \textit{spin} diskrit yang merepresentasikan momen magnetik. Meskipun pada awalnya hanya diselesaikan untuk kasus satu dimensi, solusi eksak untuk kasus dua dimensi yang dikembangkan oleh Lars Onsager kemudian berhasil membuktikan keandalan model ini dalam menjelaskan fenomena transisi fase pada suhu di atas nol mutlak. Keberhasilan matematis ini menjadikannya sebagai salah satu pilar fundamental dalam perkembangan teori fisika statistik modern yang diaplikasikan secara luas dalam berbagai disiplin ilmu.

Seiring dengan pesatnya perkembangan ilmu komputer, para peneliti mulai menyadari bahwa struktur energi pada model Ising memiliki kemiripan topologi dengan berbagai masalah optimasi kombinasional yang sangat kompleks di dunia nyata. Lanskap energi yang dibentuk oleh interaksi antar \textit{spin} dapat digunakan secara efektif untuk memetakan ruang solusi dari masalah-masalah yang tergolong dalam kelas \textit{NP-hard}, termasuk pemilihan portofolio finansial. Hal ini memicu lahirnya paradigma \textit{Ising Computing}, di mana pencarian solusi terbaik dilakukan dengan mencari konfigurasi energi terendah atau \textit{ground state} dari sistem fisik yang sedang disimulasikan melalui parameter medan magnet eksternal $h_i$. Saat ini, model Ising telah menjadi bahasa standar yang menghubungkan masalah optimasi dunia nyata dengan arsitektur komputer kuantum mutakhir berbasis \textit{Quantum Annealing} maupun algoritma variasi. Penggunaan model fisik ini memberikan perspektif baru dalam menangani kompleksitas ruang solusi yang tumbuh secara eksponensial pada masalah-masalah berskala besar.

\subsection{Representasi Fisik Aset Finansial}
Dalam bidang ekonofisika, para agen ekonomi dianalogikan sebagai sistem partikel diskrit yang saling berinteraksi secara dinamis guna mencapai titik kesetimbangan energi yang paling stabil di tengah fluktuasi pasar. Setiap aset dalam portofolio investasi kemudian direpresentasikan sebagai entitas fisik \textit{spin} ($s_i$) yang memiliki dua arah keadaan tetap di dalam ruang Hilbert yang terdefinisi secara matematis. Berdasarkan transformasi variabel yang digunakan, keadaan $-1$ merepresentasikan posisi beli (\textit{long}), sementara keadaan $+1$ digunakan untuk merepresentasikan posisi jual (\textit{short}) atau pengabaian aset dalam strategi alokasi modal. Representasi fisik ini memungkinkan analisis mendalam terhadap munculnya tren makroskopik pasar yang didasarkan pada akumulasi keputusan mikroskopik individu yang terekam dalam bentuk konfigurasi energi sistem. Interaksi antar \textit{spin} tersebut ditentukan oleh kekuatan korelasi yang dimiliki masing-masing aset melalui data ekspektasi imbal hasil dan matriks kovarians yang telah diperoleh.

Perilaku kolektif yang terjadi di pasar finansial, seperti fenomena korelasi harga antar instrumen, dapat dianalogikan sebagai interaksi magnetik antar partikel dalam sebuah kisi kristal yang terstruktur. Secara fisik, hubungan korelasi ini setara dengan parameter \textit{coupling} ($J_{ij}$), di mana kekuatan interaksi tersebut menentukan kecenderungan setiap \textit{spin} untuk menyelaraskan arah orientasinya terhadap tetangganya. Sementara itu, ekspektasi imbal hasil dari setiap aset berperan sebagai medan magnet eksternal (\textit{external field}) yang memberikan dorongan kuat pada arah orientasi \textit{spin} tertentu guna memaksimalkan utilitas investasi. Hamiltonian sistem kemudian digunakan untuk menghitung total energi yang merepresentasikan total konflik atau ketidakstabilan strategi investasi, di mana nilai energi minimum menjadi indikator utama dari solusi portofolio yang paling optimal. Integrasi antara parameter ekonomi dan variabel fisik ini menciptakan sebuah ekosistem simulasi yang mampu merefleksikan dinamika pasar finansial secara akurat.

\subsection{Pemetaan Kuantum: QUBO ke \textit{Pauli-Z}}
Untuk menjalankan proses optimasi pada arsitektur komputer kuantum, variabel biner klasik harus dipromosikan menjadi operator mekanika kuantum yang dapat diproses oleh \textit{qubit} di dalam sirkuit kuantum terpadu. Variabel keputusan biner $x_i$ dipetakan ke dalam operator \textit{Pauli-Z} ($\hat{Z}_i$) melalui sebuah transformasi linear yang bersifat eksplisit untuk menjaga integritas logika dari masalah optimasi yang dihadapi. Nilai eigen dari operator ini, yaitu $+1$ dan $-1$, secara alami mewakili status basis \textit{qubit} $|0\rangle$ dan $|1\rangle$, sehingga memungkinkan sinkronisasi yang sempurna antara data finansial dan status fisik perangkat keras. Penurunan parameter akhir dilakukan melalui substitusi variabel $x_i = (1 - \sigma_i)/2$ yang menghasilkan rumusan eksak bagi medan lokal ($h_i$) serta kekuatan interaksi ($J_{ij}$) yang akan diimplementasikan pada sirkuit kuantum sebagai berikut:
\begin{equation}
\hat{H} = \sum_{i} h_i \hat{Z}_i + \sum_{i < j} J_{ij} \hat{Z}_i \hat{Z}_j
\end{equation}
Pemetaan ini menjamin bahwa setiap interaksi kovarians antar aset dapat diterjemahkan menjadi korelasi fase antar \textit{qubit} yang saling terikat secara kuantum di dalam sirkuit optimasi.

Dengan parameter Hamiltonian yang telah berhasil dipetakan secara akurat, algoritma \textit{Variational Quantum Eigensolver} (VQE) kemudian digunakan untuk mencari solusi energi terendah dari sistem yang telah didefinisikan tersebut. VQE bekerja dengan cara mengoptimalkan parameter sirkuit secara iteratif melalui kontrol klasik guna meminimalkan nilai ekspektasi energi dari Hamiltonian Ising yang merepresentasikan risiko portofolio. Penggunaan inisialisasi cerdas dari kerangka kerja \textit{Exact Potential Game} memastikan bahwa sirkuit VQE memulai proses pencarian dari koordinat ruang Hilbert yang sudah mendekati titik solusi global yang optimal. Keberhasilan transformasi dan implementasi ini menjamin bahwa setiap hasil pengukuran akhir pada \textit{qubit} memiliki interpretasi ekonomi yang sah, logis, dan memberikan keunggulan kompetitif dalam pengambilan keputusan investasi. Dengan demikian, integrasi antara teori permainan dan komputasi kuantum memberikan solusi yang lebih presisi dalam menghadapi tantangan optimasi finansial yang bersifat \textit{NP-hard}.
